
\textbf{Existence Validated.} Behavioral fingerprints---class-specific accuracy patterns---exist in model zoos. 67.6\% of variance is behavioral, not explained by overall accuracy. This confirms that models with the same aggregate accuracy exhibit different failure patterns across classes.

\textbf{Mechanism Failed.} Simple weight statistics (25 features) do not capture this signal. Contributing factors include:

\begin{enumerate}
\item \textbf{Sample size:} 193 models versus approximately 30,000 used by Unterthiner et al. for scalar accuracy prediction. With 25 features and only 39 test samples, overfitting on training data and poor generalization are expected.
\end{enumerate}

\begin{enumerate}
\item \textbf{Feature expressivity:} Per-layer scalar statistics (mean, std, min, max, norm) may not capture the distributional or cross-layer structure encoding behavioral differences.
\end{enumerate}

\begin{enumerate}
\item \textbf{Model convergence:} Mean accuracy of 18.7\% indicates many models are undertrained. Weight statistics may be noisier for undertrained models.
\end{enumerate}

\begin{enumerate}
\item \textbf{Class 4 anomaly:} The extreme negative $R^2$ on deer (-2.74) suggests the weight features capture noise rather than signal for this low-variance class (variance = 0.013).
\end{enumerate}

\textbf{Implications.} The failure of simple statistics combined with confirmed existence of behavioral variance motivates learned representations. Neural Functionals (NF-Layers), SANE embeddings, and behavioral autoencoders with explicit behavioral supervision represent promising directions.

\textbf{Limitations.}
\begin{itemize}
\item Sample size is small (193 models versus ~30,000 in full zoo)
\item Single feature set tested (per-layer scalar statistics)
\item Single architecture family (3-conv + 2-FC CNNs on CIFAR-10)
\item No cross-architecture or cross-dataset validation
\end{itemize}
